\section*{6.1 Extended formulations and separation complexity (author source p. 13:8)}
Let us first make a short detour into extended formulations of convex polyhedra.
A polyhedron $P\subseteq\RR^r$ is a (possibly unbounded) set defined by a finite number of linear constraints. Following [20,18,2], define the extension complexity of $P$, $\xc(P)$, as the smallest $s$ such that $P$ is a linear projection of a polyhedron $Q\subseteq\RR^m$ where $Q$ can be defined using $s$ inequalities (and any number of equalities). Observe that $P$ with extension complexity $s$ can be expressed in the standard form
\[x\in P\text{ iff }\exists y\in\RR^s\ Cx+Dy=b,\ y\ge0,\]
where $C\in\RR^{t\times r}$, $D\in\RR^{t\times s}$ and $b\in\RR^t$ for some $t$.
Let $V$ be a finite subset of $\RR^r$. Given $A\subseteq V$, its separation complexity, $\sep_V(A)$, is the minimum $\xc(P)$ over all polyhedra $P\subseteq\RR^r$ with\footnote{If no such polyhedron exists, which may happen if $V$ is not convexly independent, we set $\sep_V(A):=\infty$.}
\[P\cap V=A;\]
such a $P$ is called a separating polyhedron for $A$. In other words, $\sep_V(A)$ is the smallest $s$ so that we can distinguish points in $A$ from points in $V\setminus A$ by means of a linear program with $s$ inequalities. Moreover, such a program can be rewritten as
\[x\in A\text{ iff }(x\in V\text{ and }\exists y\in\RR^s\ Cx+Dy=b,\ y\ge0).\]
The notion of separation complexity has been studied in [7,8,9] in the case when $V=\{0,1\}^n$ is the Boolean cube. The following theorem is of independent interest and can be seen as an extension of similar results in [7,9]. The proof is a considerable simplification of the previous ones.
